\section{Arithmetic packets and coefficient transport}
\label{sec:foundations}

We work with sequences, formal power series, and finite recurrence states.
Residue selectors isolate coefficients in a chosen congruence class, while
triangular transport expresses the same object in new coordinates. A valuation
filtration determines how many coordinates are needed at a given precision.
Together, these operations provide a common framework for the Euler
congruences and the finite lifting constructions developed later.

The original arithmetic seed notation
\[
\mathcal S=(A,a_0,\mathcal O,\mathcal R,\{P_{N,j}\},w,\tau_p)
\]
records a coefficient ring or module, a seed, allowed operators, relations,
residue selectors, a depth weight, and a topology. This notation describes the
ingredients of the construction. A category of seed presentations and a
notion of equivalence still require separate definitions. We therefore state
each theorem for a concrete object. Section~\ref{sec:open} discusses the
additional structure needed for the original native conjectures.

The Lean declarations cited in this section have namespace prefix
\lean{KomorebiPaperIX}. Their exact modules and source hashes are recorded in
the accompanying statement map. All belong to the preserved baseline
formalization. This section brings their statements and proofs together.

\subsection{Residue selectors and exact reconstruction}

Let $R$ be a commutative ring, $a:\N\to R$, $N\geq1$, and $0\leq j<N$.
Define
\[
(P_{N,j}a)(n)=
\begin{cases}
a(n),&n\equiv j\pmod N,\\
0,&\text{otherwise}.
\end{cases}
\]
This definition keeps the original index set. The compressed subsequence
$a(Nq+j)$, used below for recurrence packets, is a different representation of
the selected coefficients.

\begin{theorem}[Residue reconstruction]
\label{thm:residue-reconstruction}
The selectors satisfy
\[
P_{N,j}^2=P_{N,j},\qquad P_{N,i}P_{N,j}=0\quad(i\ne j),
\qquad \sum_{j=0}^{N-1}P_{N,j}=I.
\]
If $N\mid M$ and $M>0$, then
\[
P_{N,j}a=\sum_{\substack{0\leq r<M\\r\equiv j\pmod N}}P_{M,r}a.
\]
The same identities hold coefficientwise for formal power series.
\end{theorem}

\begin{proof}
At each index $n$, exactly one standard residue label equals $n\bmod N$.
Consequently a repeated selector changes nothing, two different selectors have
zero composite, and the sum retains $a(n)$ exactly once. If $N\mid M$, reducing
$n\bmod M$ modulo $N$ gives $n\bmod N$. This proves the refinement identity
at every coefficient. Equality of all coefficients proves the formal-series
versions.
\end{proof}

\begingroup\small\raggedright
\noindent\emph{Formal correspondence.}
\lean{ResiduePackets.packet_idempotent},
\lean{ResiduePackets.packet_orthogonal},
\lean{ResiduePackets.packet_reconstruction}, and
\lean{ResiduePackets.packet_refinement}. The series variants have the
prefix \lean{series_packet} in the same module.
The algebraic statements need no convergence hypothesis.
\par\endgroup

\subsection{Formal exponential coordinates and Hurwitz multiplication}

Over a characteristic-zero field $K$, put
\[
\mathcal E(a)=\sum_{n\geq0}\frac{a_n}{n!}t^n,
\qquad \mathcal D(A)_n=n![t^n]A.
\]
These inverse maps between sequences and $K[[t]]$ provide \emph{formal}
exponential coordinates. No analytic convergence is needed for the coefficient
identities below.

\begin{theorem}[Constructed unit-triangular coordinates]
\label{thm:egf-coordinates}
Let $h:\N\to K$ satisfy $h_0=1$, write $H=\mathcal E(h)$ and
$H^{-1}=\mathcal E(k)$, and define
\[
U_Ha=\mathcal D(H\mathcal E(a)).
\]
Then $U_H$ is an equivalence, with inverse
$U_H^{-1}a=\mathcal D(H^{-1}\mathcal E(a))$. Its coefficient formulas are
\[
(U_Ha)_n=\sum_{r=0}^{n}\binom nr h_r a_{n-r},\qquad
(U_H^{-1}a)_n=\sum_{r=0}^{n}\binom nr k_r a_{n-r}.
\]
The resulting triangular matrix has diagonal entries $1$. The transported
selectors
\[
\widehat P_{N,j}=U_HP_{N,j}U_H^{-1}
\]
are idempotent and pairwise orthogonal, and their sum is the identity.
\end{theorem}

\begin{proof}
The factorials are nonzero in $K$, so encoding and decoding cancel
coefficientwise. The constant coefficient of $H$ is $1$. Hence its formal
power-series inverse exists. Multiplication by $H$ and $H^{-1}$ gives the two
inverse laws. The ordinary product coefficient formula, multiplied by $n!$,
gives the displayed binomial sums. In row $n$, the coefficient of $a_n$ is
$h_0=1$ and coefficients of $a_q$ with $q>n$ vanish. Conjugating the identities
of Theorem~\ref{thm:residue-reconstruction} by the constructed equivalence proves
the transported projector identities.
\end{proof}

\begingroup\small\raggedright
\noindent\emph{Formal correspondence.}
\lean{EGF.decode_encode}, \lean{EGF.encode_decode},
\lean{EGF.egf_product}, \lean{EGF.coordinateEquiv},
\lean{EGF.coordinate_forward}, \lean{EGF.coordinate_inverse},
and \lean{EGF.transported_reconstruction}, together with the idempotence and
orthogonality declarations in \lean{EGFCoordinates}.
\par\endgroup

In characteristic $p$, factorial division is unavailable in general, but the
coefficient operation remains meaningful over every commutative ring:
\[
(f\star g)_n=\sum_{r=0}^{n}\binom nr f_r g_{n-r},\qquad
(D^d f)_n=f_{n+d},\qquad \mathbf1_n=\begin{cases}1,&n=0,\\0,&n>0.\end{cases}
\]
Here $\star$ is Hurwitz multiplication. The symbol $D^d$ denotes a coefficient
shift, which agrees with formal differentiation after characteristic-zero
encoding.

\begin{theorem}[Hurwitz product rules]
\label{thm:hurwitz-product}
For $d\geq0$,
\[
D^d(f\star g)=\sum_{r=0}^{d}\binom dr(D^r f)\star(D^{d-r}g).
\]
If $R$ has prime characteristic $p$, then
\[
D^p(f\star g)=f\star D^p g+D^p f\star g.
\]
Over a field, if $g_0=1$, a sequence $b$ is constructed recursively with
$b\star g=\mathbf1$. If also $D^p g=Dg$ in characteristic $p$, then
$D^p b=Db$.
\end{theorem}

\begin{proof}
Pascal's identity proves the first product rule, and induction proves the
higher rule. The intermediate binomial coefficients $\binom pr$ vanish in
characteristic $p$. For the inverse, solve row $n$ using the coefficient of
$b_n$, which is $g_0=1$:
\[
b_n=\mathbf1_n-\sum_{r<n}\binom nr b_r g_{n-r}.
\]
Subtract the first and $p$th product rules for $b\star g=\mathbf1$.
The assumed identity for $g$ leaves
$(D^p b-Db)\star g=0$. A second triangular induction shows that convolution
with a sequence having constant coefficient $1$ has trivial kernel.
\end{proof}

\begingroup\small\raggedright
\noindent\emph{Formal correspondence.}
\lean{Hurwitz.shift_convolution}, \lean{Hurwitz.prime_shift_convolution},
\lean{Hurwitz.convolution_left_eq_zero},
\lean{HurwitzInverse.inverse_convolution}, and
\lean{HurwitzInverse.inverse_prime_difference}.
Working directly with coefficient sequences avoids the undefined expression
$1/n!$ in $\F_p$ when $p\mid n!$.
\par\endgroup

\subsection{Triangular inversion and a finite precision filtration}

\begin{theorem}[Forward substitution and denominator clearing]
\label{thm:triangular-foundation}
Let $T_{nq}$ be coefficients over a field with $T_{nn}\ne0$ for every $n$.
For any sequence $b$, define
\[
a_n=\frac{b_n-\sum_{q<n}T_{nq}a_q}{T_{nn}}.
\]
This constructs the unique solution of
$\sum_{q=0}^{n}T_{nq}a_q=b_n$ for all $n$. If $T_{nq}$ and $b_n$ are integers
and the solution is taken over $\Q$, then
\[
\left(\prod_{r=0}^{n}T_{rr}\right)a_n\in\Z.
\]
\end{theorem}

\begin{proof}
The recursion uses only smaller indices, and substitution verifies each row.
If a second solution agrees through row $n-1$, subtracting row $n$ and
cancelling $T_{nn}$ proves equality at $n$. For integrality, use induction and
promote every earlier clearing product to $\prod_{r<n}T_{rr}$. Multiplication
by the additional factor $T_{nn}$ cancels the only new division.
\end{proof}

\begingroup\small\raggedright
\noindent\emph{Formal correspondence.}
\lean{Triangular.solve_row}, \lean{Triangular.solve_unique}, and
\lean{Triangular.solve_denominator_integral}. The clearing product bounds the
denominator. The unscaled solution can still be nonintegral.
\par\endgroup

To isolate the filtration argument, let $p$ be any element of a commutative
ring and let $w:\N\to\N\cup\{\infty\}$. Set $p^{[r]}=p^r$ for finite $r$
and $p^{[\infty]}=0$. The latter notation means the coefficient vanishes.

\begin{theorem}[Filtered finite reconstruction]
\label{thm:filtered-foundation}
For every $n,L\in\N$ and sequences $c,S$ in that ring,
\[
\begin{split}
&\sum_{r=0}^{n}\binom nr c_rp^{[w(r)]}S_{n-r}\\
&\hspace{8mm}\equiv
\sum_{\substack{0\leq r\leq n\\w(r)<L}}
\binom nr c_rp^{[w(r)]}S_{n-r}\pmod{p^L}.
\end{split}
\]
Here congruence means divisibility of the difference by $p^L$ in the
coefficient ring. In $\Zp$, it means membership in $p^L\Zp$.
\end{theorem}

\begin{proof}
The difference of the two sums consists exactly of terms with $w(r)\geq L$.
Every such term is divisible by $p^L$, including the zero terms with infinite
weight. Their finite sum has the same divisibility.
\end{proof}

\begingroup\small\raggedright
\noindent\emph{Formal correspondence.}
\lean{Filtered.filtered_triangular_reconstruction} and
\lean{Filtered.padic_filtered_reconstruction}.
For weights $2r$, retaining $r\leq M$ gives the bound $p^{2M+2}$ through
\lean{Filtered.even_weight_tail}. A sharper kernel weight can improve this
bound, as in the classical dyadic reconstruction later in the chapter. Each
coefficient row is finite, so no hypothesis $w(r)\to\infty$ is needed here.
\par\endgroup

\subsection{Finite Mahler coordinates and weighted differences}

For $f:\N\to R$, define the step-$h$ forward difference and its iterates by
\[
\Delta_h f(n)=f(n+h)-f(n),\qquad \Delta_h^0 f=f,
\]
and put $c_r(f)=\Delta_1^r f(0)$. Thus the coefficients are determined by the
differences of $f$.

\begin{theorem}[Finite Newton--Mahler identities]
\label{thm:mahler-foundation}
For every $n,s,h\in\N$,
\[
\begin{aligned}
f(n)&=\sum_{r=0}^{n}c_r(f)\binom nr,\\
c_r(\Delta_1^s f)&=c_{s+r}(f),\\
\Delta_1^s f(n)&=\sum_{j=0}^{n}c_{s+j}(f)\binom nj,\\
\Delta_h^s f(n)&=\sum_{j=0}^{s}(-1)^{s-j}\binom sj f(n+jh).
\end{aligned}
\]
The finite interpolant
$I_Nf(n)=\sum_{r=0}^{N}c_r(f)\binom nr$ agrees with $f$ for $n\leq N$.
It has the same order-$s$ unit-step difference at $n$ whenever $n+s\leq N$.
For arbitrary coefficients $c_0,\ldots,c_N$, the corresponding binomial
interpolant has exactly those coefficients and zero coefficients above $N$.
\end{theorem}

\begin{proof}
Write $E_h f(n)=f(n+h)$ for the shift operator. The unit shift is
$I+\Delta_1$. Expanding its $n$th power and evaluating
at zero gives reconstruction. Composing difference operators gives the
coefficient shift and applying reconstruction to $\Delta_1^s f$ gives the
third identity. The binomial expansion of $(E_h-I)^s$ gives the explicit
stencil. Binomial coefficients vanish when their lower index exceeds $n$,
which proves the interpolation claim. Equality of every point of the stencil
proves equality of differences. Repeated finite differences of $\binom nr$
at zero give the exact coefficient extraction.
\end{proof}

\begingroup\small\raggedright
\noindent\emph{Formal correspondence.}
\lean{Mahler.finite_mahler_reconstruction},
\lean{Mahler.coefficient_difference},
\lean{Mahler.difference_from_coefficients},
\lean{Mahler.difference_explicit},
\lean{Mahler.finite_sample_difference}, and
\lean{Mahler.coefficient_interpolant}.
Finite samples determine the stencils contained in the sample interval.
A universal tail estimate requires additional information.
\par\endgroup

For $x\in\Zp$, extend the valuation to zero by setting
\[
\operatorname{ord}_p(x)=
\begin{cases}v_p(x),&x\ne0,\\\infty,&x=0.\end{cases}
\]
Thus $L\leq\operatorname{ord}_p(x)$ is equivalent to $x\in p^L\Zp$.

\begin{theorem}[Weighted difference bound]
\label{thm:weighted-mahler-foundation}
Let $f:\N\to\Zp$ and let $w:\N\to\N\cup\{\infty\}$ satisfy
$w(s+j)\leq\operatorname{ord}_p(c_{s+j}(f))$ for $0\leq j\leq n$.
Then
\[
\operatorname{ord}_p(\Delta_1^s f(n))\geq
\min_{0\leq j\leq n}
\left\{w(s+j)+v_p\binom nj\right\}.
\]
\end{theorem}

\begin{proof}
Apply the nonarchimedean inequality to the finite coefficient formula in
Theorem~\ref{thm:mahler-foundation}. For $0\leq j\leq n$, the integer
$\binom nj$ is positive. The order of its product with $c_{s+j}(f)$ is the
sum of the two orders, including the case where that coefficient is zero.
\end{proof}

\begingroup\small\raggedright
\noindent\emph{Formal correspondence.}
\lean{PadicBounds.nat_le_order_iff},
\lean{PadicBounds.weighted_difference_bound}, and
\lean{PadicBounds.weighted_difference_divisibility}.
An expansion on all of $\Zp$ requires a convergence theorem in addition to
these finite algebraic identities and bounds.
\par\endgroup

\subsection{Unit modes and the dyadic distinction}

\begin{theorem}[Differences of unit modes]
\label{thm:unit-modes-foundation}
Over every commutative ring,
\[
\Delta_h^s(u^n)=u^n(u^h-1)^s.
\]
For any prime $p$, units $u_i\in\Zp^\times$, coefficients $a_i\in\Zp$, and
a finite index set $I$, the sequence $f(n)=\sum_{i\in I}a_i u_i^n$ satisfies
\[
\Delta_{p-1}^s f(n)\in p^s\Zp
\qquad(n,s\in\N).
\]
\end{theorem}

\begin{proof}
Induction gives the geometric identity. Reduction of a unit to $\F_p^\times$
gives $u_i^{p-1}=1$ modulo $p$. Taking $s$ powers supplies $p^s$ in each
summand. The difference operator is linear over finite sums.
\end{proof}

\begingroup\small\raggedright
\noindent\emph{Formal correspondence.}
\lean{UnitModes.geometric_difference},
\lean{UnitModes.unit_fermat_divisibility}, and
\lean{UnitModes.finite_unit_modes_divisibility}.
The result includes $p=2$, where the displayed step is $1$. Its scope here
is finite sums of modes.
\par\endgroup

\begin{theorem}[Finite unit coordinates]
\label{thm:finite-unit-foundation}
Let $p$ be prime, $L\geq0$, and
$a\in(\Z/p^{L+1}\Z)^\times$. Define
\[
\omega_L(a)=a^{p^L},\qquad \langle a\rangle_L=a\omega_L(a)^{-1}.
\]
Then
\[
a=\omega_L(a)\langle a\rangle_L,\qquad
\omega_L(a)^{p-1}=1,\qquad \langle a\rangle_L\equiv1\pmod p.
\]
The two factors are uniquely determined by these properties and are compatible
with reduction to lower precision.
For $a\in(\Z/2^{L+2}\Z)^\times$, there is instead a unique decomposition
\[
a=\epsilon(a)\langle a\rangle_4,
\qquad \epsilon(a)\in\{1,-1\},\qquad
\langle a\rangle_4\equiv1\pmod4.
\]
\end{theorem}

\begin{proof}
The unit-group exponent $p^L(p-1)$ gives the root identity for $\omega_L$.
Its residue is that of $a$, so the remaining factor has residue $1$.
If two $(p-1)$st roots have equal residue, raising them to $p^L$ makes them
equal modulo $p^{L+1}$, while each root is fixed by this exponent. This proves
uniqueness and then compatibility under reduction. In the dyadic construction,
an odd residue is exactly $1$ or $-1$ modulo $4$. Choose that sign and multiply
$a$ by it to obtain the principal factor. The two signs remain distinct
modulo $4$, which proves uniqueness.
\end{proof}

\begingroup\small\raggedright
\noindent\emph{Formal correspondence.}
\lean{FiniteUnits.finite_unit_decomposition},
\lean{FiniteUnits.decomposition_unique},
\lean{UnitRefinement.omega_compatible},
\lean{UnitRefinement.principal_compatible}, and
\lean{DyadicUnits.dyadic_unit_decomposition} with
\lean{DyadicUnits.decomposition_unique}.
The finite $p=2$ version of the first construction has trivial root-of-unity
factor. The separate modulo-$4$ construction retains the sign that reduction
modulo $2$ cannot detect.
\par\endgroup

\subsection{Recurrence packets and finite arithmetic synchronization}

\begin{theorem}[Companion jumps and packet recurrences]
\label{thm:companion-foundation}
Let $C$ be a $d\times d$ matrix over a commutative ring and
$x_n=C^n x_0$. For $N,j,k\in\N$, put $b_k=x_{Nk+j}$. Then
\[
b_k=C^j(C^N)^kx_0,\qquad b_{k+1}=C^Nb_k.
\]
If $\chi_{C^N}(X)=X^d+\sum_{r<d}q_rX^r$, then
\[
b_{k+d}=-\sum_{r=0}^{d-1}q_r b_{k+r}.
\]
Every coordinate and every linear scalar observable satisfies the same
recurrence. No invertibility or diagonalization hypothesis is required.
\end{theorem}

\begin{proof}
The jump identity follows from the power law for one matrix. Apply the
Cayley--Hamilton identity $\chi_{C^N}(C^N)=0$ to the state $b_k$.
The leading coefficient is $1$, so the final state can be isolated. A linear
observable preserves the resulting finite sum.
\end{proof}

\begingroup\small\raggedright
\noindent\emph{Formal correspondence.}
\lean{Companion.companion_jump},
\lean{Companion.packet_charpoly_recurrence}, and
\lean{Companion.packet_observable_recurrence}.
The source-sequence versions derive the packet recurrence from the relation
$x_{n+1}=Cx_n$.
For the Fibonacci matrix $C=\left(\begin{smallmatrix}1&1\\1&0\end{smallmatrix}\right)$,
its jump trace is the Lucas number $L_N$ and its determinant is $(-1)^N$.
Thus
\[
F_{N(k+2)+j}=L_NF_{N(k+1)+j}-(-1)^NF_{Nk+j}.
\]
This specialization is
\lean{FibonacciPackets.fibonacci_packet_recurrence}. It is derived from the
Fibonacci state.
\par\endgroup

\begin{theorem}[Dyadic--odd packet factorization]
\label{thm:crt-foundation}
Let $a\geq0$ and let $M$ be a positive odd integer. The Chinese remainder map
is a ring isomorphism
\[
\Z/(2^aM)\Z\simeq\Z/2^a\Z\times\Z/M\Z.
\]
For $0\leq r<2^aM$ and any coefficient sequence $f$,
\[
P_{2^aM,r}f
=P_{2^a,r\bmod 2^a}\,P_{M,r\bmod M}f.
\]
The same identity holds for formal-series coefficient selectors.
\end{theorem}

\begin{proof}
The two moduli are coprime. An index has the prescribed combined residue
exactly when it has both component residues. Applying this equivalence to
the indicator defining each selector proves the formula.
\end{proof}

\begingroup\small\raggedright
\noindent\emph{Formal correspondence.}
\lean{CRTPackets.addressEquiv}, \lean{CRTPackets.residue_equivalence},
\lean{CRTPackets.packet_factorization}, and
\lean{CRTPackets.series_packet_factorization}.
This theorem synchronizes finite arithmetic addresses. The passage from a
family of $p$-adic objects to a global integer object is a separate question,
discussed in Section~\ref{subsec:native133}.
\par\endgroup

\subsection{Polyphase notation retained from the recovered manuscript}
\label{subsec:polyphase-exposition}

The recovered Annals text develops a broader polyphase vocabulary. We explain
it through elementary identities with explicit algebraic derivations. These
identities have not been separately formalized in Lean for this consolidation.
For a sequence $a$ and $M\geq1$, let $a_r^{[M]}(q)=a_{Mq+r}$ for $0\leq r<M$.
If $A(z)=\sum_n a_nz^n$ and $A_r(w)=\sum_q a_{Mq+r}w^q$ are ordinary generating
functions, then unique quotient and remainder give
\[
A(z)=\sum_{r=0}^{M-1}z^rA_r(z^M),\qquad
\bigl(a_r^{[M]}\bigr)_s^{[N]}=a_{r+Ms}^{[MN]}.
\]
For ordinary Cauchy multiplication, the phase-$r$ product has a term without a
carry and a term with one carry:
\[
C_r(w)=\sum_{u+v=r}A_u(w)B_v(w)
+w\sum_{u+v=r+M}A_u(w)B_v(w),\qquad 0\leq u,v<M.
\]
Adding the two phase addresses gives the formula. It describes ordinary
Cauchy multiplication. Hurwitz convolution has its own binomial weights.

For $M\geq1$ and $s,r,k\in\N$, define the affine Mahler kernel by
\[
\Theta_{M,s}(r,k)=[z^r](1+z)^s\bigl((1+z)^M-1\bigr)^k
=\sum_{j=0}^{k}(-1)^{k-j}\binom kj\binom{Mj+s}{r}.
\]
It vanishes unless $k\leq r\leq Mk+s$. For a finite binomial expansion
$f(n)=\sum_{r=0}^{D}c_r\binom nr$, the $k$th Newton coefficient of
$g(q)=f(Mq+s)$ is therefore
\[
c_k(g)=\sum_{r=k}^{\min(D,Mk+s)}c_r\Theta_{M,s}(r,k).
\]
Indeed, insert the finite stencil for $\Delta_1^k g(0)$ and interchange its
two finite sums. Coefficient extraction of the binomial expansion of
$((1+z)^M-1)^k$ gives the kernel. In the binary zero-address case,
\[
\Theta_{2,0}(k+j,k)=2^{k-j}\binom kj\qquad(0\leq j\leq k),
\]
since $(2z+z^2)^k=z^k(2+z)^k$.
Exact valuations for the general kernel remain to be determined. The
recovered infinite analytic extensions and stronger historical identities
lie beyond these finite derivations. Their status is discussed in
Section~\ref{subsec:historical-scope}.
